Read in Gauge Theory Gravity, the connection and curvature forms below are the rotation-gauge field and the field strength of the Schwarzschild field in the static orthonormal frame, and the bivector they give on the plane $e_{c} \wedge e_{d}$ is $R(\hb(e_{c} \wedge e_{d})) = -\mathcal{R}(e_{c} \wedge e_{d})$ (Fact \claimref{synthesis.imported-gtg-dictionary}, clauses (ii) and (iv)). The script derives the curvature two-forms on the $(t, r, \vartheta)$ block from the connection by the second structure equations $R^{a}{}_{b} = d\omega^{a}{}_{b} + \omega^{a}{}_{c}\wedge\omega^{c}{}_{b}$, with $f'$ and $f''$ obtained by differentiation in $\Q(F)$, and finds $R^{0}{}_{1} = -(f''/2)\,\theta^{0}\wedge\theta^{1}$, $R^{0}{}_{2} = R^{1}{}_{2} = -(f'/2r)\,\theta^{0}\wedge\theta^{2}, \theta^{1}\wedge\theta^{2}$, with $-f''/2 = 2M/r^{3}$ and $-f'/2r = -M/r^{3}$; the $\hat3$ planes mirror $\hat2$ and the angular coefficient $(1 - f)/r^{2}$ is taken from Fact \claimref{pairing.imported-schwarzschild-connection-curvature}. It then lowers the index, forms the bivectors and computes the six classes; the loop statements are bilinearity, and the Weyl reading of them is Lemma \claimref{pairing.loop-classes-are-weyl-components}. That $R(\ell, e_{2})$ is a null rotation about $\ell$ is the identity $(M/r^{3})(e_{02} + e_{12}) = (M/r^{3})\,\ell e_{2}$, checked as multivectors.
